% ECONOMICS_PROBLEM: {"id": "O33-286", "title": "Aggregate productivity and distributional effects of generative AI", "jel_code": "O33", "source_id": "OP-0020"}
% !TeX program = xelatex
\documentclass[11pt,a4paper]{article}
\usepackage[margin=23mm]{geometry}
\usepackage{fontspec}
\setmainfont{Latin Modern Roman}
\usepackage{amsmath,amssymb,mathtools}
\usepackage{xcolor,enumitem,booktabs,longtable,array,xurl,hyperref}
\definecolor{muted}{HTML}{53636C}
\hypersetup{colorlinks=true,urlcolor=blue,linkcolor=blue}
\setlength{\parindent}{0pt}
\setlength{\parskip}{5pt}
\newcommand{\R}{\mathbb R}
\newcommand{\E}{\mathbb E}
\newcommand{\N}{\mathbb N}
\newcommand{\ind}{\mathbf 1}
\newcommand{\Var}{\operatorname{Var}}
\newcommand{\Cov}{\operatorname{Cov}}
\newcommand{\supp}{\operatorname{supp}}
\DeclareMathOperator*{\argmin}{arg\,min}
\DeclareMathOperator*{\argmax}{arg\,max}
\newcommand{\entrymeta}[3]{{\sffamily\small\color{muted}\textbf{\texttt{#1}}\quad|\quad #2\par #3\par}}
\newcommand{\Problemref}[1]{\texttt{#1}}
\newcommand{\JELref}[2]{\texttt{#2} (JEL \texttt{#1})}
\begin{document}
\section*{Aggregate productivity and distributional effects of generative AI}
\textbf{Problem ID:} \texttt{O33-286}\quad \textbf{JEL:} \texttt{O33}\par
% BEGIN ECONOMICS STATEMENT
\label{problem:OP-0020}
\label{atlas:prob:A10}

\noindent\textbf{Original atlas ID:} A10. \textbf{Formulation:} editorial specification of a research family.

\subsection*{Definitions and assumptions}

Fix a country, baseline date, finite horizon $H$, and a technologically explicit counterfactual excluding the specified generative-AI capabilities while holding other exogenous shocks constant. Let $d=(d_{ft})_{f,t}$ be a firm-adoption and task-use path, including complementary capital and organizational adjustment. A general-equilibrium model $\gamma$ maps that path and the common shocks into quantity output $Y_t(d;\gamma)>0$, hours $L_t(d;\gamma)>0$, wages, profits, prices, and household utilities, with an explicit equilibrium selector. Let $d^1$ denote the adoption counterfactual and $d^0$ the baseline. Define the horizon-specific labor-productivity effect
\[
\Delta_H^{\rm prod}(\gamma)=\log\frac{Y_H(d^1;\gamma)}{L_H(d^1;\gamma)}-
 \log\frac{Y_H(d^0;\gamma)}{L_H(d^0;\gamma)}.
\]
Output indices use the same accounting and quality-adjustment conventions in both paths.

\subsection*{Formal research question}

Identify or bound $\Delta_H^{\rm prod}$ and distributional changes in real labor income, profit income, employment, and household welfare, accounting for endogenous diffusion, new tasks, substitution, and market reallocation. If $\Gamma_I$ is the set of models consistent with causal workplace evidence, adoption measurement, production-network observations, and declared restrictions, report
\[
\mathcal T_I(H)=\{(\Delta_H^{\rm prod}(\gamma),\Delta_H^{\rm income}(\gamma),
 \Delta_H^{\rm welfare}(\gamma)):\gamma\in\Gamma_I\}.
\]
For fixed population groups $g$, define $\Delta_{H,g}^{\rm income}=\mathbb E[I_{g,H}(d^1;\gamma)-I_{g,H}(d^0;\gamma)]$, with real labor and profit income $I_{g,H}$ assigned by fixed ownership weights. Define $\Delta_{H,g}^{\rm welfare}$ as the corresponding difference in expected lifetime utility under fixed cardinal utility functions; the displayed vectors collect these group effects. Combine adoption telemetry with employer--employee panels and output prices to test transport of micro effects; make long-horizon projections conditional on specified innovation and diffusion paths rather than treating them as measured outcomes.

\subsection*{Interpretation and scope}

Within-firm productivity gains, economy-wide labor productivity, GDP growth, and welfare are distinct targets. General-equilibrium wage and price adjustments, business stealing, consumer gains, and new firms can change the aggregate effect relative to a sum of treatment effects. A survey association between adoption and performance does not identify a causal productivity effect without addressing selection. Task exposure is also not identical to realized adoption or displacement. The no-AI comparison must say which capabilities, complementary investments, and subsequent innovations change. Distributional gains depend on capital ownership and transfers as well as wages; a national labor share alone cannot identify who benefits. Unobserved future horizons require explicit model-based uncertainty.

\subsection*{Source audit and status}

Acemoglu--Restrepo (2018) is ambiguous because multiple relevant works exist: the published task-model paper is [S1] and the explicitly AI-focused NBER paper is [S5]. Brynjolfsson--Li--Raymond (2023) is verified as NBER Working Paper 31161; the 2025 Quarterly Journal of Economics publication is [S4]. Bonney et al. (2026) is verified as NBER 35141 and Census CES 26-25, with the six-author title-page record below; it measures diffusion and associations, not causal aggregate gains. Original link [20] identifies the earlier workplace working paper. These sources support an ongoing empirical agenda; the formal counterfactual vector above is editorial and is not a certified long-run AI conjecture.

\subsection*{References}

\begin{enumerate}

\item[S1] Daron Acemoglu and Pascual Restrepo (2018). \emph{The Race between Man and Machine: Implications of Technology for Growth, Factor Shares, and Employment}. American Economic Review 108(6), 1488--1542. \url{https://www.aeaweb.org/articles?id=10.1257/aer.20160696}. \textit{Locator:} abstract; task-based technology model. \textit{Establishes:} Growth and distributional effects of automation and new tasks in a formal model. \textit{Does not establish:} a causal long-run forecast for present generative AI.

\item[S2] Erik Brynjolfsson, Danielle Li, and Lindsey R. Raymond (2023). \emph{Generative AI at Work}. NBER Working Paper 31161. \url{https://www.nber.org/papers/w31161}. \textit{Locator:} abstract; workplace study. \textit{Establishes:} Workplace productivity and heterogeneous experience effects. \textit{Does not establish:} a 2023 journal publication or long-run economy-wide effect.

\item[S3] Kathryn Bonney, Cory L. Breaux, Emin Dinlersoz, Lucia Foster, John Haltiwanger, and Aditya Pande (2026). \emph{The Microstructure of AI Diffusion: Evidence from Firms, Business Functions, and Worker Tasks}. NBER Working Paper 35141; U.S. Census Center for Economic Studies Working Paper 26-25. \url{https://www.nber.org/papers/w35141}. \textit{Locator:} Census PDF title page, abstract, and introduction. \textit{Establishes:} Survey measurement of AI diffusion and associations with firm outcomes. \textit{Does not establish:} a causal aggregate effect or a peer-reviewed long-run equilibrium forecast.

\item[S4] Erik Brynjolfsson, Danielle Li, and Lindsey R. Raymond (2025). \emph{Generative AI at Work}. Quarterly Journal of Economics 140(2), 889--942. \url{https://doi.org/10.1093/qje/qjae044}. \textit{Locator:} abstract; firm-level deployment analysis. \textit{Establishes:} Published evidence on AI assistance in one work setting. \textit{Does not establish:} aggregate productivity, national labor shares, or long-run welfare effects.

\item[S5] Daron Acemoglu and Pascual Restrepo (2018). \emph{Artificial Intelligence, Automation and Work}. NBER Working Paper 24196. \url{https://www.nber.org/papers/w24196}. \textit{Locator:} abstract. \textit{Establishes:} A task-based framework explicitly discussing AI, displacement, and new tasks. \textit{Does not establish:} a long-run empirical general-equilibrium estimate.

\end{enumerate}

\subsection*{Common conventions: Source mathematical conventions}

Unless an entry states otherwise, agent and firm sets are finite. State and action spaces are measurable, strategies and outcome maps are measurable, probability laws are specified, and expectations exist for the targets under discussion. Infinite-horizon discount factors lie in $(0,1)$ and discounted objectives are finite. Norms, tolerances and complexity input sizes must be specified for computational comparisons. Constants or primitives that appear locally are inputs, not empirical facts asserted by this catalogue.

\textbf{Identification.} A statistical model $m\in\mathcal M$ implies an observable law $P_m$ and target $\theta(m)$. At law $P$, its identified set is
\[
\Theta(P;\mathcal M)=\{\theta(m):m\in\mathcal M,\ P_m=P\}.
\]
Point identification holds when this set is a singleton, or equivalently when $P_{m_1}=P_{m_2}$ implies $\theta(m_1)=\theta(m_2)$ for all compatible models. Sharp bounds mean bounds attained, or approached when the boundary is not attained, by compatible models. Writing this set is a definition, not a solution: the substantive task is to establish informative restrictions and derive or estimate the set. An unrestricted-confounding model may give only trivial bounds.

\textbf{Inference.} Identification, estimability and valid uncertainty quantification are separate questions. Fix $\alpha\in(0,1)$. A confidence procedure $C_n$ over a declared law class $\mathcal P$ has asymptotic uniform coverage at level $1-\alpha$ when
\[
\liminf_{n\to\infty}\inf_{P\in\mathcal P}\Pr_P\{\theta(P)\in C_n\}\geq1-\alpha.
\]
For set-identified targets the coverage event must be defined separately, for example coverage of the full identified set. If an entry uses identification-indexed models instead of uniquely indexed laws, the same coverage convention is applied over those models. Pointwise limits do not establish uniform validity, and asymptotic validity does not guarantee finite-sample precision. Sample sizes, dependence structures and weak-identification sequences are part of each application.

\textbf{Counterfactuals.} $Y_i(a)$ is a potential outcome under a specified intervention, including its timing, implementation and versions. It is not $Y_i$ conditional on voluntarily choosing $a$. A full intervention vector permits interference; shorthand scalar treatment notation does not silently impose no interference. Assignment, selection, attrition, anticipation and measurement must be specified. A claim about an instrument's local effect is not automatically a population policy effect.

\textbf{Economic equilibrium.} A counterfactual model includes best responses, constraints, feasibility, market clearing, state transitions and expectations consistent with the stated information. When several equilibria exist, either a declared selector is an input or the target is set-valued. Comparisons across policy regimes must use the stated selection convention. Reduced-form relationships are not presumed invariant after an equilibrium-changing intervention.

\textbf{Optimization and welfare.} Optimization is over the stated feasible set. If attainment is not established, $\sup$ or $\inf$ and $\varepsilon$-optimal choices replace an asserted maximizer or minimizer. For example, with value $V=\sup_{a\in\mathcal A}W(a)<\infty$, an $\varepsilon$-optimal set is $\{a\in\mathcal A:W(a)\geq V-\varepsilon\}$ for $\varepsilon>0$. Benefits, costs, risk and health or environmental outcomes can be aggregated only using declared units and conversion weights. Interpersonal utility weights, the treatment of fiscal transfers and foreign profits, discounting, and the behavioral welfare criterion are explicit inputs. A comparative-static derivative additionally requires existence and appropriate differentiability of the selected counterfactual.

\textbf{Sources.} Local labels [S1], [S2], and so forth restart in every question. Locators identify the relevant abstract, model, section or result at the level actually checked; a bibliographic or abstract-level verification is not represented as inspection of an entire proof. Working papers are distinguished from later publications and changing versions. Source summaries are paraphrased. All URL checks and editorial formulations refer to the audit date stated above.
% END ECONOMICS STATEMENT
\end{document}
